\subsection*{瞬间没有路程，怎么还能谈速度}
速度表显示某个瞬间的速度，但若把时间间隔直接设成零，路程也为零，差商无法计算。解决方法并不是允许零除以零，而是保留非零的小时间间隔，观察平均速度是否逼近一个确定值。这个稳定趋势才是瞬时速度。

令位置函数为 $f(x)=x^2$，从 $x$ 移动到 $x+h$，位移是 $2xh+h^2$。只要 $h\ne0$，平均变化率就等于 $2x+h$。在 $x=3$ 处，小正步长给略大于 $6$ 的斜率，小负步长给略小于 $6$ 的斜率，两边向同一个数靠近。

\subsection*{切线是变化率的图像}
在图上标出两点并连线，得到割线。让第二点靠近第一点，割线方向逐渐稳定为切线方向。动画中的切线不是凭目测贴在曲线上的装饰，而是差商极限的可见表示。把 $h$ 从两侧接近零，可以检查左右是否一致。

再换成 $|x|$ 的尖角。零点右侧差商为 $1$，左侧为 $-1$，没有单一斜率能同时描述两边。图形没有断开，所以连续；却没有一个统一的局部线性变化率，所以不可导。这个例子把连续和可导的区别具体化了。

学习导数时要保留两句话：先计算非零间隔下的比值，再讨论其趋势。后面看见微小量记号时，也应能回到这条解释，而不把形式上的约分误读为对零进行除法。
